\chapter{中間解析・多重性・用量反応}\label{ch:14}

\begin{goalbox}
\begin{itemize}
\item 繰り返し検定による第1種の過誤の増大を説明できる．
\item スコア統計量の部分和の独立増分構造から，情報時間 $t$ の標準化統計量の相関 $\Corr[Z(t_1),Z(t_2)]=\sqrt{t_1/t_2}$ を導ける．
\item $\alpha$ 消費関数から各解析で消費する $\alpha$ と中間解析の棄却限界値を求め，2変量正規分布を用いて最終解析の限界値を定める方法を説明できる．
\item FWER（強い制御・弱い制御）を定義し，Bonferroni 法・Holm 法を適用できる．
\item 一元配置分散分析の $F$ 検定を述べ，閉検定手順の考え方から「分散分析で保護した対比較」が FWER を強く制御する条件としない条件を説明できる．
\item 用量反応の線形対比の係数を，中心化した用量から作れる．
\end{itemize}
\end{goalbox}

\begin{exambox}
\begin{itemize}
\item \kakomon{2017}{2}：ログランク検定の群逐次デザイン．$\alpha(t)=0.05t^2$，$t=0.4$ で $\alpha=0.008$，$z_1=2.41$，
      $\Cov[Z(0.4),Z(1)]=\sqrt{0.4}$，$z_2$ を $\int_{-\infty}^{z_1}\int_{z_2}^\infty\phi_2(x_1,x_2;\sqrt{0.4})\,\dd x_2\,\dd x_1=0.042$ で定める方法，解析の追加（$t=0.7$）．
      $B(t)=\sqrt tZ(t)$ の平均と共分散は問題文で与えられた．
\item \kakomon{2024}{4}：プラセボと3用量．〔2〕Bonferroni，〔3〕分散分析をゲートとした $z$ 検定で FWER が制御される条件，〔5〕直線的用量反応の対比 $(-3,-1,1,3)/\sqrt{20}$．
      〔1〕独立な検定の FWER，〔4〕多対1比較の検定統計量の相関，〔6〕対比検定どうしの相関は，式が与えられた誘導問題である（\cref{app:F}）．
\item \kakomon{2025}{4}：2つの主要評価項目．〔1-2〕Bonferroni 法の FWER．
      〔1-1〕〔1-3〕独立な $p$ 値の FWER，〔1-4〕大域帰無仮説に対する別の棄却規則，〔2〕相関のある検定統計量の FWER は誘導問題（\cref{app:F}）．
\item \kakomon{2017}{3}：二段階デザインの $\alpha,\beta$，期待症例数．手順と二項分布表がすべて与えられた誘導問題（\cref{app:F}）．
\item \kakomon{2016}{2}：Cochran--Armitage 傾向性検定（\cref{sec:04-trend}）．
\end{itemize}
\end{exambox}

\section{問題設定}\label{sec:14-setting}
1つの試験の中で複数回・複数の仮説を検定すると，「どれか1つでも誤って有意になる」確率が名目の $\alpha$ を超える．
臨床試験では，(a) 時間的な繰り返し（中間解析），(b) 複数の群・用量，(c) 複数の評価項目，(d) 部分集団，で多重性が生じる．
本章では，これらを制御する標準的な方法を導く．

\section{群逐次デザインと \texorpdfstring{$\alpha$}{α} 消費関数}\label{sec:14-gsd}
\Hinshutsu
\subsection{情報時間と検定統計量の相関}\label{subsec:14-info}
\paragraph{情報時間}
最終解析で得られる情報量を $I_{\max}$，ある解析時点までの情報量を $I$ として，\emphx{情報時間}を $t=I/I_{\max}\in(0,1]$ とする．
ログランク検定では情報量（$U$ の分散 $V$）が近似的にイベント数に比例する（\cref{prop:12-schoenfeld}の $V\approx Dp(1-p)$）ので，
中間解析までのイベント数 $L$，最終解析の目標イベント数 $D$ に対して $t=L/D$ である．

\paragraph{スコアの部分和の構造}
情報時間 $t$ までのデータによるスコア統計量を $S(t)$ とする（ログランク検定なら，それまでのイベント時点についての $U=\sum_j(d_{1j}-E_{1j})$）．
$H_0$ の下で次が成り立つ．
\begin{enumerate}[label=(\roman*)]
\item $\E[S(t)]=0$，$\V[S(t)]\approx I(t)=t\,I_{\max}$（スコアの分散は情報量）．
\item $s<t$ のとき，増分 $S(t)-S(s)$（時点 $s$ 以降に加わるデータの寄与）は $S(s)$ と無相関．
      独立な観測の和であれば増分は過去と独立である．ログランク統計量では各時点の項 $d_{1j}-E_{1j}$ の過去を条件とした期待値が0なので（\cref{sec:09-logrank}），
      反復期待値により過去の項との共分散が0になる（情報量を確率変数でなく定数とみなす近似の下で）．
\end{enumerate}
すると $s<t$ のとき
\[
\Cov[S(s),S(t)]=\V[S(s)]+\Cov[S(s),S(t)-S(s)]=I(s)=s\,I_{\max}.
\]
標準化統計量を $Z(t)=S(t)/\sqrt{I(t)}$，$B(t)=\sqrt t\,Z(t)=S(t)/\sqrt{I_{\max}}$ とおくと
\begin{equation}
\E[B(t)]=0,\qquad \Cov[B(s),B(t)]=\frac{\Cov[S(s),S(t)]}{I_{\max}}=\min(s,t).\label{eq:14-Bcov}
\end{equation}
さらに $S(t)$ は多数の項の和なので中心極限定理により $(B(t_1),\dots,B(t_k))$ は近似的に多変量正規分布に従う．
平均0・共分散 $\min(s,t)$ の正規過程は標準 Brown 運動にほかならないので，$B(t)$ は「近似的に標準 Brown 運動」であるといわれる．
以下の計算で使うのは，有限個の時点の $(Z(t_1),\dots,Z(t_k))$ が近似的に多変量正規分布に従い，共分散が\cref{eq:14-Bcov}から決まることだけである．
\Youchui これは $H_0$ の下での\textbf{近似}であり（ログランク検定では独立増分も近似的にしか成り立たない），試験（\kakomon{2017}{2}）では $\E[B(t)]=0$，$\Cov[B(t_1),B(t_2)]=\min(t_1,t_2)$ は問題文で与えられた．

\begin{meidai}\label[meidai]{prop:14-corr}
$t_1<t_2$ のとき $\Corr[Z(t_1),Z(t_2)]=\Cov[Z(t_1),Z(t_2)]=\sqrt{t_1/t_2}$．
\end{meidai}
\begin{proof}
$Z(t)=B(t)/\sqrt t$ で $\V[Z(t)]=t/t=1$．$\Cov[Z(t_1),Z(t_2)]=\Cov[B(t_1),B(t_2)]/\sqrt{t_1t_2}=t_1/\sqrt{t_1t_2}=\sqrt{t_1/t_2}$．
\end{proof}
直観的には，$Z(t_2)$ は $Z(t_1)$ の情報に新しい（無相関な）情報を足したもので，情報の重なりの割合が相関になる．

\paragraph{繰り返し検定の問題}
$t=0.5$ と $t=1$ の2回，それぞれ片側 $z_{0.025}=1.96$ で検定すると，全体の第1種の過誤確率は
$1-\Prob(Z(0.5)\le1.96,Z(1)\le1.96)$（相関 $\sqrt{0.5}$ の2変量正規）で，数値計算すると約 $0.042$ となり，名目の0.025を大きく超える．

\subsection{\texorpdfstring{$\alpha$}{α} 消費関数}\label{subsec:14-spending}
\begin{teigi}[$\alpha$ 消費関数]\label{def:14-spending}
$\alpha(0)=0$，$\alpha(1)=\alpha$ の非減少関数 $\alpha(t)$ を定め，情報時間 $t_k$ までに累積で $\alpha(t_k)$ の第1種の過誤確率を「消費」するように棄却限界値 $c_k$ を決める
（Lan--DeMets の方法 \cite{landemets1983}）．
\end{teigi}
\begin{enumerate}
\item 第1回：$\Prob(Z(t_1)>c_1)=\alpha(t_1)$ から $c_1=z_{\alpha(t_1)}$．
\item 第 $k$ 回：$\Prob(Z(t_1)\le c_1,\dots,Z(t_{k-1})\le c_{k-1},Z(t_k)>c_k)=\alpha(t_k)-\alpha(t_{k-1})$ を満たす $c_k$ を，多変量正規分布（相関は\cref{prop:14-corr}）の数値積分で求める．
\end{enumerate}
各回の事象「それまで棄却されず，第 $k$ 回で初めて棄却」は互いに排反なので，全体の第1種の過誤確率は増分の和 $\alpha(t_K)=\alpha$ になる．
\kakomon{2017}{2} では $\alpha(t)=0.05t^2$ で $t_1=0.4$ なら $\alpha(0.4)=0.05\times0.16=0.008$，$c_1=z_{0.008}=2.41$，最終解析で消費するのは $0.05-0.008=0.042$．
最終の限界値 $c_2$ は
\[
\Prob(Z(0.4)\le c_1,\ Z(1)>c_2)=\int_{-\infty}^{c_1}\int_{c_2}^{\infty}\phi_2\bigl(x_1,x_2;\sqrt{0.4}\bigr)\,\dd x_2\,\dd x_1=0.042
\]
を満たす値で，数値計算すると $c_2\approx1.68$（名目の $z_{0.05}=1.645$ よりわずかに大きい）．
$\alpha$ 消費関数法の利点は，解析の回数や時期をあらかじめ固定しなくてよいことで，
途中で $t=0.7$ の解析を追加しても，消費する $\alpha$ は $\alpha(0.7)-\alpha(0.4)=0.0245-0.008=0.0165$，最終は $0.05-\alpha(0.7)=0.0255$ と自動的に決まる．

\begin{table}[htbp]
\centering\small
\caption{代表的な $\alpha$ 消費関数（$\alpha$ は全体の片側有意水準）}\label{tab:14-spending}
\begin{tabular}{@{}llL{62mm}@{}}
\toprule
名称 & $\alpha(t)$ & 特徴\\
\midrule
O'Brien--Fleming 型 & $2-2\Phi\bigl(z_{\alpha/2}/\sqrt t\bigr)$ & 初期はほとんど消費しない．最終の限界値は名目値に近い\\
Pocock 型 & $\alpha\log\{1+(e-1)t\}$ & 各回ほぼ同じ限界値．早期中止しやすいが最終が厳しい\\
べき型 & $\alpha t^\rho$ & $\rho=1$ で一様，$\rho$ が大きいほど初期の消費が少ない（2017年は $\rho=2$）\\
\bottomrule
\end{tabular}
\end{table}
\Youchui O'Brien--Fleming 型の式の $z_{\alpha/2}$ は両側検定を意味しない（$\alpha$ は片側水準で，$t=1$ で $2-2\Phi(z_{\alpha/2})=\alpha$ となるように置かれている）．
どの関数でも，消費する $\alpha$ の計算は上の手順と同じである．

\section{多重性と FWER}\label{sec:14-mult}
\Hissu
\subsection{FWER}
\begin{teigi}[FWER]\label{def:14-fwer}
検定する帰無仮説の族 $\{H_1,\dots,H_m\}$ について，少なくとも1つの真の帰無仮説を誤って棄却する確率を FWER（familywise error rate）という．
どの帰無仮説が真であっても（真偽のどの組合せでも）FWER $\le\alpha$ となる制御を\textbf{強い制御}，すべての帰無仮説が真のとき（完全帰無仮説）だけ保証されるものを\textbf{弱い制御}という．
\end{teigi}
真の帰無仮説が1つもなければ，誤って棄却しうる仮説がないので FWER $=0$ である．
各検定を水準 $\alpha$ で行うと，真の帰無仮説が複数あれば FWER は一般に $\alpha$ を超える．

\subsection{Bonferroni 法}
\begin{meidai}[Bonferroni 法]\label[meidai]{prop:14-bonf}
各 $H_i$ を有意水準 $\alpha/m$ で検定すると，検定統計量の依存関係によらず FWER $\le\alpha$（強い制御）．
\end{meidai}
\begin{proof}
真の帰無仮説の添字集合を $I_0$（$|I_0|=m_0\le m$），$A_i$ を「$H_i$ を棄却」とする．Boole の不等式 $\Prob(\cup A_i)\le\sum\Prob(A_i)$ より
$\mathrm{FWER}=\Prob\bigl(\bigcup_{i\in I_0}A_i\bigr)\le\sum_{i\in I_0}\alpha/m=m_0\alpha/m\le\alpha$．
\end{proof}
$p$ 値で書けば「$p_i\le\alpha/m$ なら棄却」である．\kakomon{2024}{4}〔2〕では3つの対比較なので $\alpha/3$．

\subsection{Holm 法}
\begin{meidai}[Holm 法（ステップダウン）]\label[meidai]{prop:14-holm}
$p$ 値を小さい順に $p_{(1)}\le\dots\le p_{(m)}$，対応する仮説を $H_{(1)},\dots,H_{(m)}$ とする．
$k=1,2,\dots$ の順に $p_{(k)}\le\alpha/(m-k+1)$ を調べ，成り立つ限り $H_{(k)}$ を棄却し，初めて成り立たなかったところで止める（以降は棄却しない）．
この方法は依存関係によらず FWER を強く $\alpha$ 以下に制御し，Bonferroni 法で棄却される仮説はすべて棄却する．
\end{meidai}
\begin{proof}
真の帰無仮説の数を $m_0\ge1$ とする．真の帰無仮説のうち $p$ 値が最小のものが $k$ 番目に来たとすると，それより前には偽の帰無仮説しかないので $k-1\le m-m_0$，すなわち $m-k+1\ge m_0$．
真の帰無仮説を1つでも棄却するには，この $H_{(k)}$ を棄却しなければならず，そのためには $p_{(k)}\le\alpha/(m-k+1)\le\alpha/m_0$．
よって $\mathrm{FWER}\le\Prob(\text{ある }i\in I_0\text{ で }p_i\le\alpha/m_0)\le m_0\cdot\alpha/m_0=\alpha$．
後半は，$\alpha/(m-k+1)\ge\alpha/m$ なので，Bonferroni で棄却される $p_{(1)},\dots$ は Holm でも順に棄却されることによる．
\end{proof}

\section{一元配置分散分析}\label{sec:14-anova}
準1級の範囲なので公式だけをまとめる．$a$ 群，群 $i$ の $n_i$ 人（$N=\sum_in_i$）について
\[
y_{ij}=\mu+\alpha_i+\varepsilon_{ij},\qquad \varepsilon_{ij}\iid\Normal(0,\sigma^2),\qquad \sum_{i=1}^an_i\alpha_i=0
\]
とする．群平均を $\bar y_{i\cdot}$，全平均を $\bar y_{\cdot\cdot}$ として
\[
S_T=\sum_{i}\sum_{j}(y_{ij}-\bar y_{\cdot\cdot})^2,\qquad
S_A=\sum_{i}n_i(\bar y_{i\cdot}-\bar y_{\cdot\cdot})^2,\qquad
S_E=\sum_{i}\sum_{j}(y_{ij}-\bar y_{i\cdot})^2,
\]
\begin{equation}
S_T=S_A+S_E,\qquad
F=\frac{S_A/(a-1)}{S_E/(N-a)}\ \sim\ F(a-1,\,N-a)\quad(H_0:\alpha_1=\dots=\alpha_a=0\ \text{の下で}).\label{eq:14-anova}
\end{equation}
$F>F_{a-1,N-a}(\alpha)$（$F$ 分布の上側 $100\alpha\%$ 点）で $H_0$（全群の平均が等しい）を棄却する．
$S_E/(N-a)$ は $H_0$ の真偽によらず $\sigma^2$ の不偏推定量である．$\sigma^2$ が既知なら $S_A/\sigma^2\sim\chi^2_{a-1}$（$H_0$ の下）を用いる．
$a=2$ では $F$ は2標本 $t$ 統計量の2乗に一致する．

$H_0:\mu_0=\mu_1=\dots=\mu_J$ は，各用量対プラセボの帰無仮説 $H_{0j}:\mu_0=\mu_j$ の\textbf{共通部分} $\bigcap_jH_{0j}$ である．
この見方から，分散分析を「ゲート」にした対比較の性質を次節の閉検定手順で調べる．

\section{閉検定手順}\label{sec:14-closed}
\Hissu
\begin{teiri}[閉検定手順]\label[teiri]{thm:14-closed}
族のすべての共通部分仮説 $H_I=\bigcap_{i\in I}H_i$（$I\subset\{1,\dots,m\}$，$I\ne\emptyset$）に水準 $\alpha$ の検定を用意し，
$H_i$ を「$i$ を含むすべての $I$ について $H_I$ が棄却されたとき」に棄却すると，FWER は強い意味で $\alpha$ 以下に制御される．
\end{teiri}
\begin{proof}
真の帰無仮説の添字集合を $I_0\ne\emptyset$ とすると $H_{I_0}$ は真である．真の $H_i$ $(i\in I_0)$ が1つでも棄却されるには $H_{I_0}$ が棄却されなければならず，その確率は $\alpha$ 以下．
\end{proof}
\begin{itemize}
\item Holm 法は，各 $H_I$ を Bonferroni（$\min_{i\in I}p_i\le\alpha/|I|$）で検定する閉検定手順と同じ結果を与える．
\item \textbf{固定順序（階層的）検定}：あらかじめ順序 $H_1\to H_2\to\cdots$ を決め，前の仮説がすべて棄却されたときだけ次を水準 $\alpha$ で検定する．
      $H_I$ を「$I$ の中で最も順序の早い仮説の検定」で検定する閉手順にあたり，強い制御が成り立つ．主要評価項目→副次評価項目，非劣性→優越性（\cref{sec:13-ni}）の順の検定がこれにあたる．
\end{itemize}

\paragraph{分散分析で保護した対比較（\kakomon{2024}{4}〔3〕）}
プラセボ（$j=0$）と3用量で，一元配置分散分析（水準 $\alpha$）で $H_0:\mu_0=\mu_1=\mu_2=\mu_3$ が棄却されたときだけ，
$H_{0j}:\mu_0=\mu_j$（$j=1,2,3$）をそれぞれ水準 $\alpha$ の $z$ 検定で検定する．真の帰無仮説の数で場合分けする．
\begin{itemize}
\item \textbf{すべて真}（$\mu_0=\mu_1=\mu_2=\mu_3$）：誤った棄却には分散分析での棄却が必要で，その確率は $\alpha$．よって FWER $\le\alpha$．
\item \textbf{真がちょうど1つ}（$H_{0j}$ だけ真）：誤った棄却は $H_{0j}$ の $z$ 検定の棄却を含むので FWER $\le\Prob(H_{0j}\text{ を棄却})=\alpha$．
\item \textbf{真が1つもない}：FWER $=0$．
\item \textbf{真がちょうど2つ}（例えば $\mu_0=\mu_1=\mu_2\ne\mu_3$）：$\mu_3$ が大きく異なるか症例数が大きければ分散分析はほぼ確実に棄却され，
      その後2つの真の帰無仮説 $H_{01},H_{02}$ をそれぞれ水準 $\alpha$ で検定することになる．FWER は $\Prob(H_{01}\text{ または }H_{02}\text{ を棄却})$ に近づき，これは $\alpha$ を超えうる．
\end{itemize}
したがって，各群の人数によらず FWER $\le\alpha$ が保証されるのは「$\mu_0=\mu_1=\mu_2=\mu_3$」または「$H_{0j}$ のうち真であるものが高々1つ」の場合である．
閉検定手順の言葉でいえば，2つの仮説の共通部分 $H_{0j}\cap H_{0k}$ に対する水準 $\alpha$ の検定が欠けているため，一般には弱い制御にとどまる．
プラセボと2用量（仮説が2つ）なら，共通部分は全体の帰無仮説だけなのでこの手順は閉検定手順そのものになり，強い制御が成り立つ．

\section{用量反応と線形対比}\label{sec:14-contrast}
各群の平均 $\bar Y_j\sim\Normal(\mu_j,\sigma^2/n_j)$（$j=0,\dots,J$，独立，$\sigma^2$ 既知）について，$\sum c_j=0$ を満たす対比係数 $\bm c$ による
\begin{equation}
Z_c=\frac{\sum_jc_j\bar Y_j}{\sqrt{\sigma^2\sum_jc_j^2/n_j}}\label{eq:14-contrast}
\end{equation}
は $H_0:\sum c_j\mu_j=0$ の検定統計量で，$H_0$ の下で $\Normal(0,1)$ に従う．$\sum c_j=0$ なので全群の平均が等しければ $H_0$ が成り立つ．
\begin{meidai}[直線的な用量反応の対比]\label[meidai]{prop:14-linear}
各群同数（$n_j=n$）で，真の平均が用量 $x_j$ の1次式 $\mu_j=a+bx_j$ のとき，$\E[Z_c]$ の絶対値を最大にする対比は中心化した用量 $c_j\propto x_j-\bar x$（$\bar x=\sum_jx_j/(J+1)$）である．
\end{meidai}
\begin{proof}
$\sum c_j=0$ より $\sum c_j\mu_j=b\sum c_jx_j=b\sum c_j(x_j-\bar x)$．よって $\E[Z_c]=b\sqrt n\sum c_j(x_j-\bar x)/(\sigma\|\bm c\|)$．
Cauchy--Schwarz の不等式 $|\sum c_j(x_j-\bar x)|\le\|\bm c\|\,\|\bm x-\bar x\bm1\|$ の等号は $\bm c\propto\bm x-\bar x\bm1$ のときで，これは $\sum c_j=0$ を満たす．
\end{proof}
\begin{itemize}
\item 等間隔の用量（例：0, 40, 80, 120 mg）では $x_j-\bar x=-60,-20,20,60\propto(-3,-1,1,3)$．
      ノルムを1に正規化すると $(-3,-1,1,3)/\sqrt{20}$（\kakomon{2024}{4}〔5〕）．
\item プラセボ対全用量の平均の比較は $(-3,1,1,1)$．
\item 2値データでの直線対比の検定が Cochran--Armitage 傾向性検定（\cref{sec:04-trend}）である．
\end{itemize}
\Youchui 直線対比が有意でも，用量反応が直線であることが示されたわけではない（単調でない形でも $\sum c_j\mu_j\ne0$ になりうる）．

\section{仮定}\label{sec:14-assump}
\begin{itemize}
\item 群逐次：$H_0$ の下で統計量の部分和が近似的に独立増分・正規（ログランク・スコア統計量で近似的に成立），解析時期が結果に依存しない．
\item Bonferroni・Holm：依存関係によらず有効．
\item 分散分析・対比：独立性，正規性，等分散．
\end{itemize}

\section{結果の解釈}\label{sec:14-interp}
\begin{itemize}
\item 中間解析で早期に有効中止した試験は，効果を過大推定する傾向がある（早く境界を越えたのは偶然大きく出たため）．
\item 多重性の調整で有意にならなかった副次項目の名目上の $p<0.05$ は，探索的な結果として扱う．
\item 用量反応の対比が有意でも，どの用量が有効かは各用量対プラセボの比較（多重性を調整して）で別に評価する．
\end{itemize}

\section{手法選択}\label{sec:14-choice}
\begin{itemize}
\item 中間解析：$\alpha$ 消費関数で各解析の限界値を決める．
\item 複数の仮説（複数の用量・評価項目）のどれか1つでも主張する：Bonferroni 法，より強力にするなら Holm 法．あらかじめ順序が決まっていれば固定順序検定．
\item 3群以上の平均の一様性：一元配置分散分析．分散分析をゲートにした対比較は，仮説が3つ以上なら強い制御にならない．
\item 用量反応の傾向：中心化した用量を係数とする線形対比（2値なら Cochran--Armitage）．
\end{itemize}

\section{よくある誤り}\label{sec:14-err}
\begin{warnbox}
\begin{itemize}
\item 中間解析ごとに名目の $\alpha$ で検定する．
\item $\alpha$ 消費関数で，第 $k$ 回に消費する $\alpha$ を $\alpha(t_k)$ とする（正しくは増分 $\alpha(t_k)-\alpha(t_{k-1})$）．
\item $\Corr[Z(t_1),Z(t_2)]$ を $t_1/t_2$ とする（正しくは $\sqrt{t_1/t_2}$）．$B(t)$ と $Z(t)$ の共分散を取り違える．
\item Holm 法で，一度棄却できなかった後の仮説も検定し続ける．
\item 分散分析をゲートにすれば常に FWER が強く制御されると考える．
\item 対比係数の和を0にし忘れる（全群が等しくても $H_0$ が成り立たなくなる）．
\end{itemize}
\end{warnbox}

\section{数理統計との接続}\label{sec:14-link}
\begin{linkbox}
\begin{itemize}
\item $B(t)=\sqrt tZ(t)$ の共分散構造は，独立な観測の部分和 $S_k=\sum_{i\le k}X_i$ の共分散 $\Cov(S_k,S_l)=\min(k,l)\sigma^2$ と同じである．
\item $p$ 値が帰無仮説の下で $\Unif(0,1)$ に従うこと（『現代数理統計学の基礎』定理7.22）から $\Prob(p_i\le\alpha/m)=\alpha/m$ となり，Bonferroni 法の証明が成り立つ．
\item 一元配置分散分析の $F$ 分布は，$S_A/\sigma^2$ と $S_E/\sigma^2$ が独立な $\chi^2$ 分布に従うこと（Cochran の定理）による．
\end{itemize}
\end{linkbox}

\section{例題}\label{sec:14-ex}
\begin{reidai}[$\alpha$ 消費関数]\label{reidai:14-1}
片側 $\alpha=0.025$，$\alpha(t)=0.025t^2$ で，情報時間 $t=0.5$ で中間解析，$t=1$ で最終解析を行う．
\begin{enumerate}[label=(\arabic*)]
\item 中間解析と最終解析で消費する $\alpha$ を求め，中間解析の棄却限界値を求めよ（$z_{0.00625}=2.50$）．
\item $\Corr[Z(0.5),Z(1)]$ を求め，最終解析の限界値 $c_2$ を定める式を書け（数値計算では $c_2\approx2.02$）．
\item 計画を変更して $t=0.75$ にも中間解析を追加した．2回目の中間解析と最終解析で消費する $\alpha$ を求めよ．
\item 両解析を名目の $1.96$ で行った場合の全体の第1種の過誤確率は約 $0.042$ である．このことの意味を述べよ．
\end{enumerate}
\end{reidai}

\begin{reidai}[多重性]\label{reidai:14-3}
\begin{enumerate}[label=(\arabic*)]
\item 3つの仮説の $p$ 値が $0.010,\ 0.020,\ 0.030$ であった．$\alpha=0.05$ として，Bonferroni 法と Holm 法で棄却される仮説を答えよ．
\item プラセボと2用量の3群で，一元配置分散分析（水準 $\alpha$）で全体の帰無仮説が棄却されたときだけ，各用量対プラセボを水準 $\alpha$ で検定する．
      この手順は FWER を強く制御するか．プラセボと3用量の4群ではどうか．
\end{enumerate}
\end{reidai}

\begin{reidai}[対比と多重性]\label{reidai:14-4}
プラセボと3用量（0, 40, 80, 120 mg）の各群20人で，平均が $10,12,13,15$，$\sigma^2=16$（既知）であった．
片側 $\alpha=0.025$ とし，$z_{0.025}=1.960$，$z_{0.0125}=2.241$，$z_{0.025/3}=2.394$ を用いよ．
\begin{enumerate}[label=(\arabic*)]
\item 直線的な用量反応を検出する対比係数を作り，その検定統計量を求めて判定せよ．
\item 各用量対プラセボの $Z_j=(\bar Y_j-\bar Y_0)/\sqrt{\sigma^2(1/20+1/20)}$ を求め，Bonferroni 法と Holm 法で判定せよ．
\end{enumerate}
\end{reidai}

\section{解答}\label{sec:14-sol}
\begin{kaitou}{reidai:14-1}
(1) $\alpha(0.5)=0.025\times0.25=0.00625$，最終で $0.025-0.00625=0.01875$．$c_1=z_{0.00625}=2.50$．\\
(2) \cref{prop:14-corr}より $\sqrt{0.5/1}=0.707$．$\Prob(Z(0.5)\le2.50,\ Z(1)>c_2)=0.01875$，すなわち
$\int_{-\infty}^{2.50}\int_{c_2}^{\infty}\phi_2(x_1,x_2;0.707)\,\dd x_2\,\dd x_1=0.01875$ を満たす $c_2$．\\
(3) $\alpha(0.75)=0.025\times0.5625=0.0140625$．2回目は $0.0140625-0.00625=0.0078125$，最終は $0.025-0.0140625=0.0109375$．\\
(4) 名目の限界値で2回検定すると，第1種の過誤確率が0.025から約0.042へ1.7倍程度に増える．中間解析では限界値を調整しなければならない．
\end{kaitou}

\begin{kaitou}{reidai:14-3}
(1) Bonferroni（$0.05/3=0.0167$）：$0.010$ のみ棄却．
Holm：$0.010\le0.05/3$ で棄却，$0.020\le0.05/2=0.025$ で棄却，$0.030\le0.05$ で棄却．3つとも棄却．\\
(2) 仮説は $H_{01}:\mu_0=\mu_1$，$H_{02}:\mu_0=\mu_2$ の2つで，共通部分 $H_{01}\cap H_{02}$ は全体の帰無仮説 $\mu_0=\mu_1=\mu_2$ である．
分散分析はその水準 $\alpha$ の検定，各 $z$ 検定は $H_{0j}$ の水準 $\alpha$ の検定なので，この手順は閉検定手順（\cref{thm:14-closed}）であり強い制御が成り立つ
（場合分けでも，真の帰無仮説が2つ（全部）・1つ・0の各場合に FWER $\le\alpha$）．
4群では，真の帰無仮説がちょうど2つの場合（例：$\mu_0=\mu_1=\mu_2\ne\mu_3$）に FWER が $\alpha$ を超えうるので，強い制御にならない（\cref{sec:14-closed}）．
\end{kaitou}

\begin{kaitou}{reidai:14-4}
(1) 中心化した用量 $-60,-20,20,60$ より $\bm c=(-3,-1,1,3)$（\cref{prop:14-linear}）．
$\sum c_j\bar Y_j=-30-12+13+45=16$，$\sqrt{\sigma^2\sum c_j^2/n}=\sqrt{16\times20/20}=4$，$Z_c=4.0>1.960$ で有意（増加傾向あり）．\\
(2) $\SE=\sqrt{16\times2/20}=\sqrt{1.6}=1.265$．$Z_1=2/1.265=1.58$，$Z_2=3/1.265=2.37$，$Z_3=5/1.265=3.95$．
Bonferroni（限界値 $2.394$）：$Z_3$ のみ有意．
Holm：最大の $Z_3=3.95>2.394$ で棄却，次の $Z_2=2.37>z_{0.025/2}=2.241$ で棄却，$Z_1=1.58<1.960$ で止める．$H_{03},H_{02}$ を棄却．
\end{kaitou}

\begin{summarybox}
\begin{itemize}
\item 情報時間 $t=I/I_{\max}$（ログランクなら $L/D$）．スコアの部分和が $H_0$ の下で平均0・分散 $\propto$ 情報量・無相関な増分をもつことから $\Cov[B(s),B(t)]=\min(s,t)$（近似），$\Corr[Z(t_1),Z(t_2)]=\sqrt{t_1/t_2}$．
\item $\alpha$ 消費：第1回 $c_1=z_{\alpha(t_1)}$，第 $k$ 回は増分 $\alpha(t_k)-\alpha(t_{k-1})$ を多変量正規で．2回なら $\Prob(Z_1\le c_1,Z_2>c_2)=\alpha(1)-\alpha(t_1)$．
\item FWER の強い制御：Bonferroni（$\alpha/m$），Holm（$p_{(k)}\le\alpha/(m-k+1)$ のステップダウン），閉検定手順，固定順序検定．
\item 一元配置分散分析：$S_T=S_A+S_E$，$F=\{S_A/(a-1)\}/\{S_E/(N-a)\}\sim F(a-1,N-a)$．
\item 分散分析で保護した対比較：仮説が2つなら閉検定手順で強い制御．3つ以上では，真の帰無仮説が2つ以上あり，かつ全部ではないとき FWER が $\alpha$ を超えうる．
\item 直線対比：中心化した用量 $(-3,-1,1,3)$，$\sum c_j=0$．
\end{itemize}
\end{summarybox}
